\input{../doc-class-cours.tex}

\begin{document}

\begin{center}
  \textit{Pourquoi apprendre alors que l’ignorance est instantanée ?} - \textbf{Bill Watterson}
\end{center}

\setlength{\columnseprule}{0pt}

\begin{multicols}{2}
\subsection*{Ex1 : Conjecture}

 \begin{enumerate}
  \item[1.] Calculer l'aire des carrés : $A_1$ et $A_2$.
  \item[2.] Calculer l'aire du carré : $A_3$.
  \item[3.] Calculer la longueur du côté du carré : $c_3$.
\end{enumerate} 

 \subsection*{Ex2 : Rédaction}
 \begin{enumerate}
  \item[2.] Le triangle MEX est rectangle en X. \newline
  On a MX = 24.2cm et ME= 35.6cm. Calculer EX.
\end{enumerate}\columnbreak

 \begin{figure}[H]
    \centering
    \includegraphics[width=0.4\linewidth]{4x3-pythagore/ex1a.pdf}
  \end{figure}  

\end{multicols}

\begin{multicols}{2}

 \subsection*{Ex3 : Probabilité}
On a une urne avec 3 balles bleues, 2 balles rouges et 5 balles vertes. 
 \begin{enumerate}
  \item[1.] Quelle est la probabilité de tirer une balle rouge ?
  \item[2.] Quelle est la probabilité de tirer une balle qui n'est pas bleue ? 
\end{enumerate}

 \subsection*{PB1 : théorique}
Soit ABCD un quadrilatère ci-contre. On a AH = 6.8cm.
 \begin{enumerate}
  \item[1.] Calculer le périmètre du quadrilatère ABCD.
  \item[2.] Calculer l'aire du quadrilatère ABCD.
\end{enumerate} \columnbreak

 \begin{figure}[H]
    \centering
    \includegraphics[width=0.3\linewidth]{4x3-pythagore/pb1.pdf}
  \end{figure}

  \end{multicols}


 \begin{figure}[H]
    \centering
    \includegraphics[width=0.3\linewidth]{4x3-pythagore/pb2.pdf}
  \end{figure}

 \subsection*{PB2 : Pont}

On souhaite consolider un pont de longueur AH = 420m. 

Le pont passe par dessus une vallée située à une altitude HJ = 280m. 

Le pont possède un mat en son milieu d'une hauteur EF = 135m. 

On relie les points B,D,E,G, I par un hauban (un câble en acier).

 \begin{enumerate}
  \item[1.] Calculer la longueur totale du hauban. 
\end{enumerate}

Pour consolider ce pont, l'ingénieur de génie civil et l'architecte du bureau d'étude préparent un devis. (le prix estimé des travaux)

\begin{itemize}[label={$\bullet$}]
  \item Câble acier : 155€ le mètre de câble.
  \item Mat en ciment avec le ferraillage : 40€ le mètre. 
\end{itemize}

 \begin{enumerate}
  \item[2.] Calculer le montant du devis. 
\end{enumerate}

\end{document}